\documentclass{article}

\title{Sales Forecasting System \\ Complete Technical Guide}
\author{DevOps, MLOps, GenAI Integration}
\date{January 2026}

\begin{document}

\maketitle

\tableofcontents

\section{Introduction}

This project demonstrates a complete modern data science and DevOps architecture combining multiple technologies:

\begin{itemize}
\item Database management with SQLite
\item Interactive analysis with Jupyter Notebook
\item Machine learning predictions
\item Generative AI integration
\item Data visualization
\item REST API development
\item Container orchestration
\item Continuous integration and deployment
\end{itemize}

\section{Part 1: SQLite Database}

\subsection{Purpose}

The SQLite database stores all historical sales data. It serves as the single source of truth for our forecasting system.

\subsection{Schema}

The sales table contains these columns:

\begin{itemize}
\item id - Primary key, auto-incremented record identifier
\item product - Product name (Widget A, Widget B, etc.)
\item region - Geographic region (North, South, East, West)
\item date - Transaction date in YYYY-MM-DD format
\item units\_sold - Integer number of units sold
\item revenue - Decimal revenue amount
\end{itemize}

\subsection{Sample Data}

Example records from the sales table:

\begin{verbatim}
id | product  | region | date       | units_sold | revenue
1  | Widget A | North  | 2025-01-01 | 10         | 100.00
2  | Widget B | South  | 2025-01-01 | 5          | 50.00
3  | Widget A | North  | 2025-01-02 | 12         | 120.00
\end{verbatim}

\subsection{Creation}

Create the database using sales.sql:

\begin{verbatim}
sqlite3 sales.db < sales.sql
\end{verbatim}

\section{Part 2: Jupyter Notebook Analysis}

\subsection{Overview}

The Jupyter notebook contains 10 interconnected cells that form the analysis pipeline.

\subsection{Cell Structure}

Cell 1: Data Loading and Export

Connects to SQLite, loads all sales data, exports to CSV for backup, filters by region.

Cell 2-3: Package Installation

Installs Plotly, Dash, and other required libraries for visualization and web framework.

Cell 4: OpenAI API Integration

Calls OpenAI's ChatGPT model with historical sales data to generate AI predictions.

Cell 5: Response Parsing

Extracts numerical predictions from AI text response using regex parsing.

Cell 6: Actual Sales Visualization

Creates interactive Plotly line chart showing historical sales trends.

Cell 7: Predicted Sales Visualization

Creates interactive Plotly chart showing forecast predictions.

Cell 8: FastAPI Endpoint Definition

Defines REST API endpoints for /sales and /predict routes.

Cell 9: API Testing

Tests endpoints and validates responses with error handling.

Cell 10: Server Launch

Starts the FastAPI server with uvicorn for local testing.

\subsection{Key Libraries}

pandas - Data manipulation and analysis

sqlite3 - Database connectivity

plotly - Interactive visualizations

scikit-learn - Machine learning algorithms

openai - OpenAI API client

fastapi - Modern web framework

uvicorn - ASGI server

\section{Part 3: Machine Learning Model}

\subsection{Algorithm}

LinearRegression from scikit-learn provides baseline forecasting by learning linear trends.

\subsection{How It Works}

The model fits a line through historical data points and extrapolates to future dates.

Example: If sales increased by 2 units per day, the model predicts this trend continues.

\subsection{Implementation}

\begin{verbatim}
from sklearn.linear_model import LinearRegression
import pandas as pd

# Convert dates to ordinal numbers
X = data['date_ordinal'].values.reshape(-1, 1)
y = data['units_sold'].values

# Train model
model = LinearRegression()
model.fit(X, y)

# Predict next 7 days
future_dates = [[date1], [date2], ..., [date7]]
predictions = model.predict(future_dates)
\end{verbatim}

\subsection{Strengths}

Simple and interpretable

Fast training and prediction

Low computational overhead

Works with limited data

\subsection{Limitations}

Cannot capture complex nonlinear patterns

Assumes constant trend

Sensitive to outliers

No seasonal adjustment

\section{Part 4: Generative AI Integration}

\subsection{OpenAI API}

Uses ChatGPT (gpt-4o-mini model) for intelligent sales analysis.

\subsection{Process}

1. Format historical sales data as text
2. Send prompt to OpenAI API
3. Receive AI-generated forecast
4. Parse numbers from response
5. Use as additional prediction source

\subsection{Implementation}

\begin{verbatim}
from openai import OpenAI

client = OpenAI(api_key='YOUR_KEY')

history = data.to_dict('records')
prompt = f'Predict sales: {history}'

response = client.chat.completions.create(
    model='gpt-4o-mini',
    messages=[{'role': 'user', 'content': prompt}]
)

prediction = response.choices[0].message.content
\end{verbatim}

\subsection{Response Parsing}

Uses regex to extract numbers from AI text response:

\begin{verbatim}
import re
numbers = re.findall(r'\b\d+\b', prediction)
forecast = [int(n) for n in numbers[:7]]
\end{verbatim}

\subsection{API Considerations}

Requires valid API key from platform.openai.com

Incurs cost per API call

gpt-4o-mini offers good balance of speed and accuracy

\section{Part 5: Data Visualization}

\subsection{Plotly Dashboard}

Interactive charts for exploring sales data and predictions.

\subsection{Chart Types}

Line charts for time-series trends

Scatter plots with markers for predictions

Interactive legends and hover tooltips

Responsive sizing for different displays

\subsection{Components}

\begin{verbatim}
import plotly.graph_objects as go

# Actual sales chart
fig1 = go.Figure()
fig1.add_trace(go.Scatter(
    x=data['date'],
    y=data['units_sold'],
    mode='lines',
    name='Actual Sales'
))

# Predicted sales chart
fig2 = go.Figure()
fig2.add_trace(go.Scatter(
    x=future_dates,
    y=predictions,
    mode='lines+markers',
    name='Forecast'
))
\end{verbatim}

\subsection{Features}

Hover information shows exact values

Click legend items to toggle visibility

Download charts as PNG

Zoom and pan capabilities

\section{Part 6: FastAPI REST API}

\subsection{Overview}

Exposes predictions via HTTP endpoints for external applications.

\subsection{Endpoints}

GET /sales - Returns all sales records as JSON

GET /predict - Returns forecast for specified region

GET /docs - Interactive API documentation (Swagger UI)

\subsection{Implementation}

\begin{verbatim}
from fastapi import FastAPI

app = FastAPI()

@app.get('/sales')
def get_sales():
    return data.to_dict('records')

@app.get('/predict')
def predict(region: str = 'North'):
    return {'region': region, 'forecast': predictions}
\end{verbatim}

\subsection{Running}

\begin{verbatim}
uvicorn api:app --host 0.0.0.0 --port 8000
\end{verbatim}

Access at http://localhost:8000/docs

\subsection{Usage}

\begin{verbatim}
curl http://localhost:8000/sales
curl http://localhost:8000/predict?region=North
\end{verbatim}

\section{Part 7: Docker Containerization}

\subsection{Purpose}

Packages application with all dependencies in isolated container.

\subsection{Dockerfile}

\begin{verbatim}
FROM python:3.11
WORKDIR /app
COPY . .
RUN pip install -r requirements.txt
CMD ["python", "app.py"]
\end{verbatim}

\subsection{Building}

\begin{verbatim}
docker build -t sales-app:latest .
\end{verbatim}

\subsection{Running}

\begin{verbatim}
docker run -d -p 8000:8000 sales-app:latest
docker logs <container-id>
docker stop <container-id>
\end{verbatim}

\subsection{Benefits}

Consistent environment across machines

Simplified deployment process

Easy version management

Isolation from host system

\section{Part 8: Kubernetes Orchestration}

\subsection{Overview}

Manages containerized applications at scale across multiple servers.

\subsection{Deployment Configuration}

\begin{verbatim}
apiVersion: apps/v1
kind: Deployment
metadata:
  name: sales-app
spec:
  replicas: 3
  selector:
    matchLabels:
      app: sales-app
  template:
    metadata:
      labels:
        app: sales-app
    spec:
      containers:
      - name: sales-app
        image: sales-app:latest
        ports:
        - containerPort: 8000
\end{verbatim}

\subsection{Service Configuration}

\begin{verbatim}
apiVersion: v1
kind: Service
metadata:
  name: sales-service
spec:
  type: LoadBalancer
  selector:
    app: sales-app
  ports:
  - port: 80
    targetPort: 8000
\end{verbatim}

\subsection{Deployment Commands}

\begin{verbatim}
kubectl apply -f deployment.yaml
kubectl apply -f service.yaml
kubectl get pods
kubectl get services
kubectl logs <pod-name>
\end{verbatim}

\subsection{Features}

Automatic pod restart on failure

Load balancing across replicas

Rolling updates with zero downtime

Horizontal scaling

\section{Part 9: Jenkins CI/CD Pipeline}

\subsection{Purpose}

Automates testing, building, and deployment when code changes.

\subsection{Jenkinsfile}

\begin{verbatim}
pipeline {
    agent any
    
    stages {
        stage('Checkout') {
            steps {
                git url: 'https://github.com/user/repo.git'
            }
        }
        
        stage('Build') {
            steps {
                sh 'docker build -t sales-app .'
            }
        }
        
        stage('Test') {
            steps {
                sh 'python -m pytest tests/'
            }
        }
        
        stage('Deploy') {
            steps {
                sh 'kubectl apply -f deployment.yaml'
            }
        }
    }
}
\end{verbatim}

\subsection{Pipeline Stages}

Checkout - Pull code from GitHub

Build - Create Docker image

Test - Run unit tests

Deploy - Update Kubernetes deployment

\subsection{Triggering}

Webhook from GitHub on push

Scheduled builds at specific times

Manual trigger via Jenkins UI

\subsection{Benefits}

Automated testing catches errors early

Consistent build process

Fast deployment cycles

Complete deployment history

\section{System Architecture}

\subsection{Data Flow}

SQLite Database stores historical sales data

Jupyter loads data and generates analysis

ML model learns patterns from historical data

OpenAI API provides intelligent predictions

Plotly creates interactive visualizations

FastAPI serves predictions via HTTP

Docker containers package the application

Kubernetes orchestrates container deployment

Jenkins automates the entire pipeline

\subsection{Integration Points}

Database provides input data

ML and AI produce predictions

Visualization and API consume predictions

Containers enable portability

Kubernetes provides scalability

Jenkins automates deployment

\section{Getting Started}

\subsection{Prerequisites}

Python 3.9 or higher

pip package manager

Docker installed

kubectl configured for Kubernetes cluster

Jenkins server running

\subsection{Installation Steps}

1. Clone repository from GitHub

2. Install Python dependencies:
\begin{verbatim}
pip install -r requirements.txt
\end{verbatim}

3. Create SQLite database:
\begin{verbatim}
sqlite3 sales.db < sales.sql
\end{verbatim}

4. Set OpenAI API key:
\begin{verbatim}
export OPENAI_API_KEY='your-key-here'
\end{verbatim}

5. Launch Jupyter notebook:
\begin{verbatim}
jupyter notebook sales.ipynb
\end{verbatim}

6. Execute cells in order from Cell 1 to Cell 10

7. Access API documentation at http://localhost:8000/docs

\section{Troubleshooting}

\subsection{Common Issues}

Database connection failed - Ensure sales.db exists and has correct permissions

API key invalid - Verify OpenAI key is correct and has remaining credits

Port already in use - Kill existing process or use different port

Docker build fails - Check Docker daemon is running and Dockerfile is valid

Kubernetes pod crashes - Check logs with kubectl logs command

\subsection{Debugging}

View database contents: sqlite3 sales.db

Check API logs: docker logs container-name

Inspect pod status: kubectl describe pod pod-name

Monitor resource usage: kubectl top nodes

\section{Summary}

This integrated system demonstrates:

- Data persistence with SQLite
- Interactive analysis with Jupyter
- Machine learning predictions
- Generative AI integration
- Data visualization with Plotly
- Web API with FastAPI
- Container management with Docker
- Orchestration with Kubernetes
- Automation with Jenkins

All components work together to create a production-ready forecasting platform capable of analyzing historical data, generating predictions, and serving results via scalable cloud infrastructure.

\end{document}
